% Generated draft: requires mathematical and source review.
\documentclass[11pt]{article}
\usepackage[margin=25mm]{geometry}
\usepackage{fontspec}
\setmainfont{Latin Modern Roman}
\usepackage{amsmath,amssymb,mathtools,unicode-math}
\setmathfont{Latin Modern Math}
\usepackage{enumitem,needspace,xurl,hyperref}
\setlength{\parindent}{0pt}
\setlength{\parskip}{5pt}
\setlength{\emergencystretch}{3em}
\newcommand{\N}{\mathbb{N}}
\newcommand{\Z}{\mathbb{Z}}
\newcommand{\Q}{\mathbb{Q}}
\newcommand{\R}{\mathbb{R}}
\newcommand{\C}{\mathbb{C}}
\newcommand{\F}{\mathbb{F}}
\newcommand{\A}{\mathbb{A}}
\newcommand{\PP}{\mathbb{P}}
\newcommand{\E}{\mathbb{E}}
\newcommand{\eps}{\varepsilon}
\newcommand{\dd}{\,\mathrm{d}}
\begin{document}
\section*{UnsolvedMath ID: 2302012}
\textbf{Source problem code:} catalogue entry
\subsection{Definitions and mathematical statement}
Let $f(z)=\sum_{n\geq0}a_nz^n$ be a nonconstant entire function of finite order
\[
 \rho(f)=\limsup_{r\to\infty}\frac{\log T(r,f)}{\log r}<\infty,
 \qquad \Lambda_f=\{n\geq0:a_n\ne0\}.
\]
Here $T$ is the Nevanlinna characteristic. For $a\in\C$, its finite-value deficiency is
\[
 \delta(a,f)=1-\limsup_{r\to\infty}\frac{N(r,a;f)}{T(r,f)},
\]
where $N(r,a;f)$ counts zeros of $f-a$ in the integrated Nevanlinna sense, with multiplicities.

The source calls $\Delta$ the \emph{maximal density} of the nonzero coefficients, but gives no definition in the supplied excerpt. This is a density of $\Lambda_f$, not of the numerical coefficient values. A possible gap-theory convention is the upper P\'olya density, expressed using short relative intervals by
\[
 \limsup_{\eta\uparrow1}\ \limsup_{R\to\infty}
 \frac{\#(\Lambda_f\cap(\eta R,R])}{(1-\eta)R}.
\]
This contextual formula must be checked against the source; ordinary upper asymptotic density or upper uniform density must not be silently substituted.

\textbf{Question.} With $\Delta$ interpreted as the intended maximal density, does
\[
 \rho(f)\Delta<\tfrac12
 \quad\Longrightarrow\quad
 \delta(a,f)<1\quad\text{for every finite }a\in\C
\]
hold for every entire function in the stated class? The assertion excludes deficiency exactly one; it does not assert absence of every positive deficiency, and it excludes the value $\infty$ from its conclusion.

The supplied summary reports openness in the 2018 anniversary edition, with ``no progress reported,'' and no post-2018 resolution located. It separately claims contributor verification of the $\Delta=0$ endpoint (attributed to Fuchs, 1963), the $\rho<1/2$, $\Delta=1$ case via the classical $\cos\pi\rho$ theorem, and a reduction from K\H{o}vari's conjectured minimum-modulus bound, Problem 2.12a. It also reports a quantitative estimate
\[
 \sum_a\delta(a,f)\leq C\lambda(f)\Delta,
 \qquad \lambda(f)=\liminf_{r\to\infty}\frac{\log T(r,f)}{\log r},
\]
attributed to Fuchs (1969) and Murai (1983), which it says yields the desired conclusion only when $\rho\Delta<1/C$ for a large absolute constant $C$.

Finally, the summary reports Mittag--Leffler extremal values $\delta=1-\sin(\pi\sigma)$ at $\rho\Delta=\sigma\in[1/2,1]$ and says all known deficiency-one examples have $\rho\Delta\geq1$. The displayed formula gives zero deficiency at $\sigma=1/2$; no strictly positive boundary deficiency is inferred from it. All these claims, including the contributor's purported verifications, remain unreviewed supplied evidence.
\subsection{Short English statement}
If an entire function has order $\rho$ and its nonzero Taylor coefficients have maximal density $\Delta$ with $\rho\Delta<1/2$, can a finite value have Nevanlinna deficiency exactly one? The density convention is missing, and the supplied partial results and claimed reductions require review.
\subsection{Sources}
\begin{itemize}
\item Local UnsolvedMath registry, permanent ID 2302012, source code catalogue entry, ``Research Problems in Function Theory --- Problem 2.12''; status partially solved, classification PARTIAL-PROGRESS, updated 2026-08-06. Its summary separately reports the main problem still open and describes contributor-claimed checks and reductions.
\item Hayman, \emph{Research Problems in Function Theory} (2018), Problem 2.12, \url{https://arxiv.org/abs/1809.07200}. The record reports source-TeX extraction and access on 2026-07-29; the 2018 update was retained for status triage. Current status and rights require review; the record describes publicly available arXiv source TeX.
\item Supplied partial-progress references: Fuchs (1963, 1969), Murai (1983), the classical $\cos\pi\rho$ theorem, K\H{o}vari's conjecture in Problem 2.12a of the same source, and unspecified Mittag--Leffler extremals. Complete bibliographic details and proofs are absent.
\end{itemize}
\end{document}
